For simple random walk $(S_j)$ on $\mathbb Z^2$, put
\[
 f_N(x)=\#\{0\le j\le N:S_j=x\},\quad
 M_N=\max_x f_N(x),\quad F_N=\{x:f_N(x)=M_N\}.
\]
The requested probabilities are
\[
 \mathbb P(|F_N|=r\text{ infinitely often})=
 \begin{cases}1,&r=3,\\0,&r\ge4.\end{cases}\tag{1}
\]
The deep external input used here is Hao--Li--Okada--Zheng, \emph{Favorite Sites for Simple Random Walk in Two and More Dimensions}, Theorem 1.1 in version 2, \url{https://arxiv.org/abs/2409.00995}, which states
\[
 \limsup_{N\to\infty}|F_N|=3\quad\hbox{almost surely}.\tag{2}
\]
This theorem includes both planar assertions. T\'oth's earlier eventual-three theorem concerns one-dimensional walk, as the introduction of the cited primary paper explicitly explains; it cannot by itself be used for the planar upper bound.

We use (2) as an attributed theorem rather than claim to reproduce its delicate estimates on near-favorite sites. The deductions below clarify the simultaneous events, prove the complementary cases $r=1,2$, and derive a pathwise bound for the number of sites that have ever been favorites.
